% Brief (tikz-diagram-craft).
% Reader question: when does observer a1 know what worker a2 knows?
% Claim: each agent has a directed vision cone V(a). a2 lies inside V(a1), so
%   a1 observes a2; phi lies inside V(a2), so a2 observes phi. Composing the
%   two observations gives a1 the higher-order fact K_{a1}K_{a2}phi. Points
%   outside every cone are epistemic shadow -- unobserved by anyone.
% Evidence: chapter-0-prerequisites.tex, "Dynamic Epistemic Logic and Big
%   Brother Logic" -- Charrier et al. 2014: "if agent a1 observes agent a2
%   (a2 in V(a1)), a1 knows what a2 observes, yielding K_{a1}K_{a2}(phi)".
%   [internal; Charrier et al. 2014]
% Must distinguish: the two agents' cones; the two observation relations
%   (a1-sees-a2, a2-sees-phi) that compose into the higher-order fact; the
%   shadow where neither cone reaches.
% Grammar: a real plane (not pgfplots -- nothing here is a computed curve):
%   two agents as named points, each with a directed wedge, phi as a plain
%   datum inside the inner wedge, a ruled legend composing the two numbered
%   observation relations into the verdict. Rejected: the v1 figure's
%   numbered coordinate grid (-7..10, -3..3) -- nothing in the chapter or the
%   caption ever cites a coordinate, so the grid was unearned precision, and
%   its three verdict cards repeated the identical formula twice and a third,
%   unrelated "invariant verification" card -- a content bug, not a styling
%   one. No robot icons: an agent is a named circle, the way an automaton's
%   state is.
% Five-second test: "a2's dot sits inside the wide cone; phi's dot sits
%   inside the narrow cone grown out of a2 -- so a1 knows a2 knows phi."
\pdfigurehue{pdcobalt}%
\begin{tikzpicture}[pd figure,x=1cm,y=1cm]
  % ---- the grid -------------------------------------------------------------
  \def\Wl{0.0}  \def\Wr{11.6}  \def\Hb{0.55}  \def\Ht{5.85}
  \coordinate (a1) at (1.35,2.35);
  \coordinate (a2) at (6.05,3.45);
  \coordinate (phi) at (10.05,4.35);
  \def\angA{13.2}  \def\spreadA{21}  \def\radA{8.1}   % V(a1): wide, long
  \def\angB{12.7}  \def\spreadB{17}  \def\radB{6.15}  % V(a2): narrower
  \def\trim{0.52}  % clears the 8.5mm agent node, so no edition's pattern
                    % fill ever needs to occlude a stroke under its own text
  % ---- the field: unobserved by default, named once -----------------------
  \draw[pd neutral fill] (\Wl,\Hb) rectangle (\Wr,\Ht);
  \node[pd label,anchor=south west,text=pdinkmuted] at (\Wl+0.18,\Hb+0.14)
    {epistemic shadow};
  % ---- V(a1): wide, reaches far enough to hold a2 --------------------------
  % the wedge's near-apex edge is trimmed clear of the agent node (see
  % \trim above) rather than meeting at the node's exact centre, so the
  % stroke never needs the node's own fill to occlude it.
  \fill[pd focus fill]
    ($(a1)+({\angA+\spreadA}:\trim)$) -- ($(a1)+({\angA+\spreadA}:\radA)$)
      arc[start angle={\angA+\spreadA},end angle={\angA-\spreadA},radius=\radA]
      -- ($(a1)+({\angA-\spreadA}:\trim)$) -- cycle;
  \node[pd focus label,anchor=south west] at ($(a1)+(1.2,1.55)$) {$V(a_1)$};
  % ---- V(a2): narrower, reaches far enough to hold phi ---------------------
  \fill[pd legible fill]
    ($(a2)+({\angB+\spreadB}:\trim)$) -- ($(a2)+({\angB+\spreadB}:\radB)$)
      arc[start angle={\angB+\spreadB},end angle={\angB-\spreadB},radius=\radB]
      -- ($(a2)+({\angB-\spreadB}:\trim)$) -- cycle;
  \node[pd legible label,anchor=south west] at ($(a2)+(0.55,1.35)$) {$V(a_2)$};
  % ---- the two agents and the objective, named nodes so every incident
  % line clips to the node's own boundary -- never through its text, in any
  % edition's fill treatment (hatched editions do not occlude an overlaid
  % line the way a solid tint does).
  \node[pd automaton focus,minimum size=8.5mm] (node-a1) at (a1) {$a_1$};
  \node[pd automaton state,minimum size=8.5mm] (node-a2) at (a2) {$a_2$};
  \node[pd datum] (node-phi) at (phi) {};
  \node[pd label,anchor=west] at ($(phi)+(0.28,0)$) {$\phi$};
  % ---- the two observation relations, each an arrow with a badge ----------
  \draw[pd focus arrow] (node-a1) -- (node-a2);
  \draw[pd legible arrow] (node-a2) -- (node-phi);
  \node[pd badge] at ($(a1)!0.58!(a2)$) {1};
  \node[pd badge] at ($(a2)!0.58!(phi)$) {2};
  % ---- the ruled legend: the two observations compose into the verdict ----
  % badges are set as bold numerals, not redrawn circles, so the table never
  % depends on font inheritance into a nested tikzpicture.
  \node[anchor=north west,text width=11.5cm,align=left] at (\Wl,\Hb-0.35) {%
    \renewcommand{\arraystretch}{1.28}%
    \pdfiglabel
    \begin{tabular}{@{}cp{3.55cm}p{4.65cm}@{}}
      \toprule
      & \textbf{condition} & \textbf{reading} \\
      \midrule
      \textbf{1} & $a_2 \in V(a_1)$ & $a_1$ observes $a_2$ \\
      \textbf{2} & $\phi \in V(a_2)$ & $a_2$ observes $\phi$ \\
      $1 \wedge 2$ & \multicolumn{2}{l}{the composed verdict: $a_1$ knows $a_2$ knows $\phi$} \\
      \bottomrule
    \end{tabular}};
\end{tikzpicture}
